\section{Theorem C: universal rows and the infinite regime hierarchy}\label{sec:C}
\textbf{Status.} The universal row recurrence (Theorem~\ref{thm:Cuniversal}), its run hierarchy (Corollary~\ref{cor:C runs}), the adjacent lower-boundary obstruction (Proposition~\ref{prop:Ckmin}), and twelve certified parameter intervals (Theorem~\ref{thm:Chierarchy}) are proved. The upper endpoint formula and a parameter-uniform symbolic proof remain \emph{Open}.

We investigate the empirical starting-position template
\begin{equation}\label{eq:template}
\mathcal{T}(k,A):\qquad P_0(n)\iff n\bmod p\in R,\qquad
p=2k+A+2,\quad R=\{0,A,k+A+1\},
\end{equation}
where $k,A$ are positive integers and $P_0(n)$ means that the unrestricted start with $n$ counters is losing. These observations concern this game's specific recurrence, not a claim previously established in the cited literature. The background on subtraction games and no-repetition variants motivates the question \cite{winningways2001,gurvich2012}.

\subsection{Uniform analytic results: proved}
\begin{proposition}[Immediate upper bound on any valid template]\label{prop:Cupper}
If $\mathcal{T}(k,A)$ holds for $D(k,1)$, then necessarily $k<A$.
\end{proposition}
\begin{proof}
Equation~\eqref{eq:template} declares $n=A$ a losing unrestricted start, since $A\in R$. If $k\ge A$, however, removing exactly $A$ counters is a legal \emph{first} move from $n=A$ and leaves an empty heap, where the next player has no legal move. Therefore the first player wins from $n=A$, a contradiction.
\end{proof}

\begin{proposition}[Arithmetic collapse of the three-residue template]\label{prop:Ccollapse}
Let $p=2k+A+2$ and $R=\{0,A,k+A+1\}\pmod p$ with $k,A\ge1$. The binary $p$-periodic indicator of $R$ has a smaller positive period if and only if $k=A-1$. In that case $p=3A$, $R=\{0,A,2A\}$, and the least positive period is $A$.
\end{proposition}
\begin{proof}
The residues satisfy $0<A<k+A+1<p$, so there are exactly three. If the indicator has a smaller period, a nonzero translation of the cyclic group $\mathbb Z/p\mathbb Z$ preserves $R$. Every orbit of a fixed translation has equal size, which must divide $|R|=3$. Since a nonzero translation has no fixed point, its orbit size must equal three. The invariant set is therefore the orbit $\{0,p/3,2p/3\}$. The smaller positive residue is $A$, so $A=p/3$, equivalently $3A=2k+A+2$, or $k=A-1$. Conversely, if $k=A-1$, then $p=3A$ and $R=\{0,A,2A\}$; membership depends only on divisibility by $A$, and $A$ is clearly least because only residue zero is selected modulo $A$.
\end{proof}

\subsection{The universal game: exact row structure}\label{sec:Cuniversal}
Let the \emph{universal game} be $D(\infty,1)$: removals satisfy $1\le r\le n$, together with the usual restriction $|r-s|>1$ after a previous move. For $n\le k$, the finite game's row is the universal row restricted to $1\le s\le k$, because the move bound does not remove any candidate move. Write $Q_n$ for the universal row of previous removals making $(n,s)$ a $\Pset$-position, and set
\[
r(n)=\begin{cases}1,&Q_n\ne\varnothing,\\0,&Q_n=\varnothing.\end{cases}
\]

\begin{theorem}[Universal row theorem]\label{thm:Cuniversal}
In the universal game, for every $n\ge5$,
\[
Q_n=\begin{cases}[n-1,\,n+1],&r(n)=1,\\ \varnothing,&r(n)=0.\end{cases}
\]
Moreover, for every $m\ge5$,
\begin{equation}\label{eq:recursion}
r(2m)=1-r(m),\qquad
r(2m+1)=1-\max\{r(m),r(m+1)\},
\end{equation}
with base values $r(0),\ldots,r(9)=1,1,1,1,1,0,1,0,1,1$.
\end{theorem}
\begin{proof}
By strong induction on $n$. The base rows are computed directly from Theorem~\ref{thm:Ainterval}:
\[
Q_0=[1,\infty),\quad Q_1=Q_2=[1,2],\quad Q_3=\{2\},\quad Q_4=\{3\},
\quad Q_5=\varnothing,\quad Q_6=[5,7],\quad Q_7=\varnothing,\quad Q_8=[7,9],\quad Q_9=[8,10],
\]
so $r(0),\ldots,r(9)=1,1,1,1,1,0,1,0,1,1$.

Assume the claim for all $5\le j<n$ with $n\ge10$. A removal $r$ belongs to $A_n$ exactly when $r\in Q_{n-r}$. The value $r=n$ always belongs to $A_n$, because $(0,n)$ is terminal and hence a $\Pset$-position.

If $n-r\ge5$, the induction hypothesis gives $r\in Q_{n-r}$ exactly when $Q_{n-r}\ne\varnothing$ and $|r-(n-r)|\le1$, that is, exactly when $r\in\{\lfloor n/2\rfloor,\lceil n/2\rceil\}$ and the corresponding row is nonempty. If $1\le n-r\le4$, then $r\ge n-4\ge6$ while $Q_1,Q_2\subseteq[1,2]$, $Q_3=\{2\}$, $Q_4=\{3\}$; hence none of these $r$ belongs to $A_n$. Thus $A_n$ consists of $n$ together with at most two midpoint values (only one candidate for even $n$).

If $n=2m$, the only possible midpoint value is $r=m$, which belongs to $A_n$ exactly when $r(m)=1$. When $r(m)=0$ we have $A_{2m}=\{2m\}$, so $Q_{2m}=[2m-1,2m+1]$ and $r(2m)=1$. When $r(m)=1$ we have $\{m,2m\}\subseteq A_{2m}$; these values differ by $m\ge5>2$, so Theorem~\ref{thm:Ainterval} gives $Q_{2m}=\varnothing$ and $r(2m)=0$. Hence $r(2m)=1-r(m)$.

If $n=2m+1$, the midpoint values are $m$ and $m+1$, belonging to $A_n$ exactly when $r(m+1)=1$ and $r(m)=1$, respectively. If both $r(m)=r(m+1)=0$, then $A_{2m+1}=\{2m+1\}$, so $Q_{2m+1}=[2m,2m+2]$ and $r(2m+1)=1$. Otherwise at least one midpoint value joins $2m+1$, and the extremal elements of $A_n$ differ by at least $m\ge5>2$, so $Q_{2m+1}=\varnothing$ and $r(2m+1)=0$. Hence $r(2m+1)=1-\max\{r(m),r(m+1)\}$.

In both nonempty cases $A_n=\{n\}$, so $Q_n=[n-1,n+1]$. This completes the induction.
\end{proof}

\begin{corollary}[Run hierarchy]\label{cor:C runs}
The maximal runs of $r$ are exactly:
\begin{itemize}\itemsep0.1em
\item $r=1$ on $[0,4]$, $[6,6]$, $[8,10]$;
\item for every odd $j\ge3$ and $\varepsilon\in\{0,1,2\}$, $r=1$ on the interval
\[
\Bigl[G_{j,\varepsilon},\ G_{j,\varepsilon}+L_{j,\varepsilon}-1\Bigr],
\]
where
\begin{align*}
G_{j,0}&=\frac{2+5\cdot2^{j}}{3}, & L_{j,0}&=\frac{2^{j-1}-1}{3},\\
G_{j,1}&=\frac{2+8\cdot2^{j}}{3}, & L_{j,1}&=\frac{2^{j+1}-1}{3},\\
G_{j,2}&=\frac{2+11\cdot2^{j}}{3}, & L_{j,2}&=\frac{5\cdot2^{j}-1}{3},
\end{align*}
\item the intervening $0$-runs have lengths
\[
\frac{5\cdot2^{j-1}+1}{3},\qquad \frac{2^{j}+1}{3},\qquad \frac{4\cdot2^{j}+1}{3}
\]
after the first, second, and third run of group $j$, respectively.
\end{itemize}
Consequently the $1$-run starts are $0,6,8$ together with $G_{j,\varepsilon}$ for all odd $j\ge3$, $\varepsilon\in\{0,1,2\}$.
\end{corollary}
\begin{proof}
We first record the doubling lemma. If $[u,v]$ is a maximal $0$-run of $r$ with $u\ge6$, then $[2u,2v]$ is a maximal $1$-run. Indeed, for $u\le m\le v$ we have $r(2m)=1$; for $u\le m<v$ we have $r(2m+1)=1-\max\{0,0\}=1$; and $r(2v+1)=1-\max\{0,r(v+1)\}=0$ because $v+1$ lies in the following $1$-run. On the left, $r(2u-1)=1-\max\{r(u-1),r(u)\}=1-\max\{1,0\}=0$.

Conversely, every maximal $1$-run $[p,q]$ with $p\ge11$ is the double of a maximal $0$-run with start at least $6$. If $p=2m+1$ is odd, then $r(p)=1$ forces $r(m)=r(m+1)=0$; the $0$-run containing $m$ and $m+1$ doubles to a maximal $1$-run containing $p$, hence equal to $[p,q]$, contradicting odd $p$. Hence $p=2u$ is even. Then $r(2u)=1-r(u)$ forces $r(u)=0$; the $0$-run $[u,v]$ containing $u$ doubles to a maximal $1$-run containing $p$, hence equal to $[p,q]$.

Thus $1$-runs above $10$ correspond bijectively to $0$-runs with start at least $6$. The three gaps following the runs of group $j$ map respectively to the third run of group $j$ and the first two runs of group $j+2$. Explicitly,
\[
2\Bigl(G_{j,0}+L_{j,0}\Bigr)=G_{j,2},\qquad
2\Bigl(G_{j,1}+L_{j,1}\Bigr)=G_{j+2,0},\qquad
2\Bigl(G_{j,2}+L_{j,2}\Bigr)=G_{j+2,1}.
\]
and the corresponding right endpoints match. Since the base runs and gaps are verified directly, induction generates exactly the claimed list.
\end{proof}

\begin{table}[h]
\centering
\caption{First runs of the universal sequence and the closed forms.}\label{tab:Cruns}
\begin{tabular}{rrrr}
\toprule
$j$ & $[G_{j,0},\,G_{j,0}+L_{j,0}-1]$ & $[G_{j,1},\,\dots]$ & $[G_{j,2},\,\dots]$\\
\midrule
3 & $[14,14]$ & $[22,26]$ & $[30,42]$\\
5 & $[54,58]$ & $[86,106]$ & $[118,170]$\\
7 & $[214,234]$ & $[342,426]$ & $[470,682]$\\
9 & $[854,938]$ & $[1366,1706]$ & $[1878,2730]$\\
11 & $[3414,3754]$ & $[5462,6826]$ & $[7510,10922]$\\
13 & $[13654,15018]$ & $[21846,27306]$ & $[30038,43690]$\\
\bottomrule
\end{tabular}
\end{table}

The formulas were additionally checked against the recursion \eqref{eq:recursion} through $n=10^{7}$: the $33$ maximal $1$-runs in that range have exactly the stated starts, and every complete run has the stated length and gap (\texttt{src/large\_scale\_checks.py}). The recursion itself was checked against the exact solver through $n=20000$, independently of the induction.

\subsection{The infinite regime hierarchy}\label{sec:Chierarchy}
For odd $j\ge3$ define
\[
A_{j,0}=\frac{2+5\cdot2^{j}}{3},\qquad
A_{j,1}=\frac{2+8\cdot2^{j}}{3},\qquad
A_{j,2}=\frac{2+11\cdot2^{j}}{3}.
\]
These are precisely the three $1$-run starts of group $j$ in Corollary~\ref{cor:C runs}.

\begin{theorem}[Certified regime hierarchy]\label{thm:Chierarchy}
For $j\in\{5,7,9,11,13\}$:
\begin{enumerate}\itemsep0.2em
\item $\mathcal{T}(k,A_{j,0})$ holds throughout
\[
\frac{2^{j+2}-2}{3}\le k\le A_{j,0}-1;
\]
\item $\mathcal{T}(k,A_{j,1})$ holds throughout
\[
\frac{11\cdot2^{j-1}-2}{3}\le k\le K_j,
\]
where $K_5=78$, $K_7=313$, $K_9=1251$, $K_{11}=5007$, $K_{13}=20025$;
\item $\mathcal{T}(A_{j,2}-1,A_{j,2})$ holds, with least starting period $A_{j,2}$. This certifies the collapse case, not the exclusion of all other $k$.
\end{enumerate}
The small group $j=3$ has certified cases $A=14$ at $k=11$ and $A=22$ at $k=17$, plus the collapse $A=30$ at $k=29$. Every parameter in the twelve listed intervals is certified for all $n$ by finite windows plus induction. Both adjacent parameters fail for each interval; this establishes maximal contiguous intervals, not global uniqueness. The separate checks through $n=10^7$ cover the endpoints of the ten intervals with $j\le11$. The $j=13$ extension uses finite-window certificates instead.
\end{theorem}
\begin{proof}[Exact finite certificates and induction]
For each pair $(k,A)$, the exact interval recurrence produces a repeated window of $k$ consecutive rows with displacement $p=2k+A+2$, checks the template through the complete bridge $t+k+p-1$, and verifies row equality on the remaining computed range. Lemmas~\ref{lem:Bwindow} and~\ref{lem:Bbridge} extend these finite premises to every $n$. The original \texttt{src/verify\_hierarchy.py} certifies $2583$ pairs; \texttt{src/verify\_j13.py} adds $7740$, for $10323$ total. Their files are \texttt{data/certificates/regime\_hierarchy\_certificates.csv} and \texttt{data/certificates/regime\_hierarchy\_j13.csv}. Adjacent counterexamples appear in Table~\ref{tab:Cboundary}. Independent solver comparisons for selected earlier parameters are recorded in \texttt{regime\_hierarchy\_summary.json}; the extension uses the same certifier. Proposition~\ref{prop:Ccollapse} supplies the least period for certified collapse cases, but does not establish that the template holds without a game certificate.
\end{proof}

\begin{table}[h]
\centering
\caption{Certified intervals and third-position collapse cases. ``Collapse certified'' does not assert exclusion of other parameters.}\label{tab:Chierarchy}
\begin{tabular}{rrrrrrr}
\toprule
$j$ & $A_{j,0}$ & $k$-interval & $A_{j,1}$ & $k$-interval & $A_{j,2}$ & status\\
\midrule
3 & 14 & $11$--$11$ & 22 & $17$--$17$ & 30 & collapse certified\\
5 & 54 & $42$--$53$ & 86 & $58$--$78$ & 118 & collapse certified\\
7 & 214 & $170$--$213$ & 342 & $234$--$313$ & 470 & collapse certified\\
9 & 854 & $682$--$853$ & 1366 & $938$--$1251$ & 1878 & collapse certified\\
11 & 3414 & $2730$--$3413$ & 5462 & $3754$--$5007$ & 7510 & collapse certified\\
13 & 13654 & $10922$--$13653$ & 21846 & $15018$--$20025$ & 30038 & collapse certified\\
\bottomrule
\end{tabular}
\end{table}

\begin{table}[h]
\centering
\caption{First counterexamples at the regime boundaries. ``Predicted'' and ``actual'' refer to whether the unrestricted start is a $\Pset$-position.}\label{tab:Cboundary}
\begin{tabular}{rrrrrr}
\toprule
$A$ & $k_{\min}-1$ & first $n$ & $k_{\max}+1$ & first $n$ & predicted/actual at $k_{\max}+1$\\
\midrule
54 & 41 & 42 & 54 & 54 & P/N\\
86 & 57 & 58 & 79 & 492 & P/N\\
214 & 169 & 170 & 214 & 214 & P/N\\
342 & 233 & 234 & 314 & 1314 & P/N\\
854 & 681 & 682 & 854 & 854 & P/N\\
1366 & 937 & 938 & 1252 & 5238 & P/N\\
3414 & 2729 & 2730 & 3414 & 3414 & P/N\\
5462 & 3753 & 3754 & 5008 & 20942 & P/N\\
13654 & 10921 & 10922 & 13654 & 13654 & P/N\\
21846 & 15017 & 15018 & 20026 & 83746 & P/N\\
\bottomrule
\end{tabular}
\end{table}

\begin{proposition}[Lower boundary mechanism]\label{prop:Ckmin}
Let $e\ge10$ be an even right endpoint of a maximal $1$-run of the universal sequence, with $r(e/2)=0$. Then for every $A>e$ the template $\mathcal{T}(e-1,A)$ fails at $n=e$: it predicts an $\Nset$-position, but $e$ is a $\Pset$-position.
\end{proposition}
\begin{proof}
In $D(e-1,1)$, rows below $e$ are universal rows restricted to previous removals at most $e-1$, and the removal $r=e$ is illegal. A removal $r<e$ belongs to $A_e$ exactly when $r\in Q_{e-r}$; restriction does not affect this test since $r\le e-1$. If $e-r\ge5$, Theorem~\ref{thm:Cuniversal} forces $|2r-e|\le1$; if $1\le e-r\le4$, then $r\ge e-4\ge6$ while $Q_{e-r}\subseteq[1,4]$. Since $e$ is even and $Q_{e/2}=\varnothing$ by hypothesis, $A_e=\varnothing$. Hence the start $n=e$ is a $\Pset$-position. On the other hand $e\bmod p=e$ with $p=2(e-1)+A+2>e$ and $e\notin\{0,A,e+A\}$, so $\mathcal{T}(e-1,A)$ predicts an $\Nset$-position.
\end{proof}

For the runs above $10$, each right endpoint is twice the right endpoint of its preimage $0$-run under the doubling lemma. Thus Proposition~\ref{prop:Ckmin} explains the adjacent lower-boundary failures for the listed intervals with $j\ge5$. It does not exclude disconnected valid intervals at smaller $k$. The exceptional $j=3$ boundaries are checked separately.

\begin{remark}[Upper boundary mechanism and open formula]\label{rem:Cupper}
For the checked groups $j\in\{7,9,11,13\}$, the first counterexample at $k=K_j+1$ is $n=p+A_{j,1}$. Here $Q$ denotes rows of the finite game: the removal $A_{j,1}/2$ belongs to $Q_{p+A_{j,1}/2}$, so removing it leads to a losing state and is a winning first move. For $j=13$, this membership is absent at $k=20025$ and present at $k=20026$, where $p=61900$ and the counterexample is $n=83746$. This membership test is not monotone in $k$ and does not justify binary search for a global endpoint. The observed deficits $d_j=A_{j,1}-K_j$ are
\[
d_3=5,\quad d_5=8,\quad d_7=29,\quad d_9=115,\quad d_{11}=455,\quad d_{13}=1821,
\]
and no closed form is known. The cases $j=3$ and $j=5$ have different failure modes ($n=A+k+1$ and $n=2p$, respectively).
\end{remark}

\subsection{Exact fixed-parameter certificate tables}
\begin{longtable}{rrrrrl}
\caption{Theorem C fixed-parameter certificates for $A=54$.}\label{tab:Ccert}\\
\toprule $k$ & candidate $p$ & least & $t$ & bridge end & SG \\
\midrule\endfirsthead\toprule $k$ & candidate $p$ & least & $t$ & bridge end & SG \\
\midrule\endhead
42 & 140 & 140 & 111 & 292 & PASS \\
43 & 142 & 142 & 117 & 301 & PASS \\
44 & 144 & 144 & 117 & 304 & PASS \\
45 & 146 & 146 & 115 & 305 & PASS \\
46 & 148 & 148 & 116 & 309 & PASS \\
47 & 150 & 150 & 112 & 308 & PASS \\
48 & 152 & 152 & 114 & 313 & PASS \\
49 & 154 & 154 & 0 & 202 & PASS \\
50 & 156 & 156 & 0 & 205 & PASS \\
51 & 158 & 158 & 0 & 208 & PASS \\
52 & 160 & 160 & 0 & 211 & PASS \\
53 & 162 & 54 & 0 & 214 & PASS \\
\bottomrule\end{longtable}

\begin{longtable}{rrrrrl}
\caption{Counterexamples to extending the fixed-$A$ formulas.}\label{tab:Cfail}\\
\toprule $k$ & $A$ & first $n$ & predicted P? & actual P? & direct P/N \\
\midrule\endfirsthead\toprule $k$ & $A$ & first $n$ & predicted P? & actual P? & direct P/N \\
\midrule\endhead
41 & 54 & 42 & no & yes & PASS \\
54 & 54 & 54 & yes & no & PASS \\
55 & 54 & 54 & yes & no & PASS \\
56 & 54 & 54 & yes & no & PASS \\
57 & 54 & 54 & yes & no & PASS \\
57 & 86 & 58 & no & yes & PASS \\
79 & 86 & 492 & yes & no & PASS \\
80 & 86 & 167 & yes & no & PASS \\
\bottomrule\end{longtable}
\paragraph{Interpretation.} The $k=54,\ldots,57$ failures of $A=54$ at $n=54$ follow already from Proposition~\ref{prop:Cupper}. The $k=79$, $A=86$ failure at $n=492$ is particularly instructive: several earlier blocks of starting outcomes happen to agree with the proposed formula, but its recurrence ultimately disagrees. One must not infer a theorem from a long numerical scan or a finite streak of residue matches.

\subsection{What remains Open}
\begin{conjecture}[Genuine/non-genuine triples]\label{conj:Chierarchy}
For every odd $j\ge3$, the three $1$-run starts of group $j$ produce exactly two genuine regimes and one collapse-only value. The certified intervals and collapse cases through $j=13$ support this conjecture. The exclusion of additional intervals or genuine parameters at the third value is not implied by these certificates or by a coarse scan.
\end{conjecture}
\statusopen\quad The following remain unproved: a parameter-uniform symbolic derivation of the transition rows and of the template on a whole regime; a closed form for the upper endpoints $K_j$ (equivalently the deficits $d_j$); a proof that no genuine regime exists at a value $A$ that is not a universal $1$-run start; and a proof of the genuine/non-genuine pattern for all $j$. The finite certificates do not by themselves yield these statements.

\subsection{Next verifiable proof obligations}
The next step is to derive a finite, symbolic, piecewise affine formula for $Q_n$ uniformly in $k$ on a regime, prove each branch stable under the exact set recurrence, and justify the allowed parameter inequalities, in particular the upper endpoint condition in Remark~\ref{rem:Cupper}. A certificate-assisted theorem for finitely many $k$ values must be distinguished from such a uniform proof.\par
